\begin{textsample}{NEG Dim 2 – persona\_gemini – Score 1.00 – t881\_gemini.txt}  \label{ex:f2_neg_009}
A global movement to combat environmental pollution is growing as governments implement bans on plastics.
France has banned thin single-use plastic \textbf{bags} as part of a broader European Union effort.
In a more comprehensive measure, Morocco has enacted a ban on the production, import, \textbf{sale}, and distribution of plastic \textbf{bags}, although it has met with resistance from stockpiling.
Action is also being taken at the state and local levels.
India 's Karnataka state has banned various plastic items, including \textbf{bags}, plates, and microbeads, and authorities have seized 39,000kg of illegal plastic.
The city of San Francisco began by banning plastic \textbf{bags} in 2007, later banned plastic water \textbf{bottles} in 2014, and has also banned styrofoam, a material that is hard to recycle.
Other examples include Honolulu, which introduced a \textbf{bag} ban in 2015, and Tasmania, which reduced its \textbf{bag} usage by 350,000 as early as 2003.
Ethiopia also implemented a ban in 2011, linking it to a green energy initiative.
These bans address the critical issue of plastic pollution in our oceans, which is harmful to marine life.
To be part of this growing solution, Greenpeace urges you to join the Plastics Pledge.

% matched lemmas: bag, bottle, sale
\end{textsample}
